\section{Introduction}
\label{sec:introduction}

Modern vision architectures have diverged rapidly in structural design: vision transformers use layer normalization and global self-attention~\cite{touvron2021training}, hybrid architectures combine depthwise separable convolutions with layer normalization~\cite{liu2022convnet}, while classical CNNs retain batch normalization~\cite{he2016deep}. Yet optimizer selection has not adapted accordingly---practitioners still default to AdamW or SGD regardless of these structural differences, treating optimizer choice as secondary to architecture and data.

This neglect is consequential. SGD collapses to $11.51\%$ accuracy on layer-normalized ConvNeXt-T while achieving $69.50\%$ on batch-normalized ResNet-50, suggesting normalization type as a critical factor---a hypothesis we examine in Section~\ref{sec:spectral_analysis}. The same task remains accessible to other optimizers, revealing that the optimizer-architecture interaction is not a minor hyperparameter issue but a fundamental compatibility constraint that current practice ignores.

Why do some optimizers transfer across architectures while others fail on certain architectures? Existing first-order methods---AdamW with diagonal curvature estimates, SGD with raw gradient direction---do not explicitly account for how architectural design reshapes the spectral properties of weight gradients. Prior optimizer comparisons typically evaluate single architectures~\cite{schmidt2021crowdedvalley} or lack controlled protocol isolation~\cite{choi2020optimizers}, making it impossible to disentangle optimizer capability from recipe engineering.

We hypothesize that optimizers exploiting \emph{weight-matrix spectral structure}---through Newton--Schulz orthogonalization or Kronecker-factored preconditioning~\cite{bernstein2024old}---should exhibit architecture-robust performance precisely because they adapt to the spectral geometry regardless of how normalization and layer composition reshape it. We test this through a controlled study of three spectral optimizers (Muon~\cite{jordan2024muon}\footnote{At the time of writing, the Muon reference is a blog post and has not been peer-reviewed.}, SOAP~\cite{vyas2024soap}, Shampoo~\cite{gupta2018shampoo}) against AdamW~\cite{loshchilov2019decoupled} and SGD, spanning three architecture families, three datasets, and two training horizons, with independently tuned learning rates per optimizer--architecture pair (Section~\ref{sec:optim_protocol}).

Our experiments confirm the family-level hypothesis across all tested configurations, revealing architecture-dependent mechanisms.

While BOCB~\citep{li2024bocb} diagnoses coupling bias and Kunstner et al.\ explain why SGD fails on transformers~\citep{kunstner2024noise}, neither tests spectral methods nor identifies the architectural mediator. Our work bridges these lines: we evaluate the spectral family that Bernstein \& Newhouse unified theoretically, and identify normalization---not architecture topology---as the mechanism that determines when spectral preconditioning helps.

\paragraph{Contributions.}
\begin{itemize}[leftmargin=*,nosep]
\item The first controlled empirical test of whether Bernstein \& Newhouse's theoretical family~\cite{bernstein2024old} shares practical properties, revealing that normalization type---not architecture family---is the primary mediator of optimizer-architecture coupling.
\item Causal evidence that the spectral advantage is mediated by normalization type: replacing BatchNorm with LayerNorm on ResNet-50 collapses AdamW by ${\sim}17$\,pp while Muon is preserved (+1.64\,pp), and layer-wise stable-rank-divergence correlations reverse sign between CNNs ($\rho{=}{+}0.45$) and transformers ($\rho{=}{-}0.40$).
\item Practical guidance: spectral methods are architecture-robust candidates across the training conditions tested, achieving superior accuracy at substantially lower training cost.
\item A corruption robustness evaluation on CIFAR-100-C showing that spectral optimizers, particularly Muon, improve not only clean accuracy but also robustness to common input corruptions (mCE 89--93\% on three of four architecture configurations vs.\ AdamW baseline of 100\%).
\end{itemize}
